% Generated draft: requires mathematical and source review.
\documentclass[11pt]{article}
\usepackage[margin=25mm]{geometry}
\usepackage{fontspec}
\setmainfont{Latin Modern Roman}
\usepackage{amsmath,amssymb,mathtools,unicode-math}
\setmathfont{Latin Modern Math}
\usepackage{enumitem,needspace,xurl,hyperref}
\setlength{\parindent}{0pt}
\setlength{\parskip}{5pt}
\setlength{\emergencystretch}{3em}
\newcommand{\N}{\mathbb{N}}
\newcommand{\Z}{\mathbb{Z}}
\newcommand{\Q}{\mathbb{Q}}
\newcommand{\R}{\mathbb{R}}
\newcommand{\C}{\mathbb{C}}
\newcommand{\F}{\mathbb{F}}
\newcommand{\A}{\mathbb{A}}
\newcommand{\PP}{\mathbb{P}}
\newcommand{\E}{\mathbb{E}}
\newcommand{\eps}{\varepsilon}
\newcommand{\dd}{\,\mathrm{d}}
\begin{document}
\section*{UnsolvedMath ID: 700425}
\textbf{Source problem code:} catalogue entry
\subsection{Definitions and mathematical statement}
A variety of groups is the class of all groups satisfying some fixed collection of group identities, each universally quantified over its variables. A subvariety of a variety \(\mathcal V\) is a variety \(\mathcal W\) with \(\mathcal W\subseteq\mathcal V\); it is proper if \(\mathcal W\ne\mathcal V\). Write \(\operatorname{Sub}(\mathcal V)\) for the set of subvarieties of \(\mathcal V\).

Define the ordinal height of \(\mathcal V\) by
\[
h(\mathcal V)=\sup\left\{\operatorname{otp}(C,\subseteq): C\subseteq\operatorname{Sub}(\mathcal V)\setminus\{\mathcal V\},\ (C,\subseteq)\text{ is well-ordered}\right\}.
\]
Here \(\operatorname{otp}\) denotes order type, and the order is increasing inclusion. Equivalently, the ordinals in this supremum are those \(\alpha\) for which there is a family \((\mathcal V_\beta)_{\beta<\alpha}\) of proper subvarieties of \(\mathcal V\) satisfying
\[
\mathcal V_\beta\subsetneq\mathcal V_\gamma\qquad(\beta<\gamma<\alpha).
\]
The empty chain is allowed and has order type \(0\); no continuity condition at limit indices is imposed.

The problem asks whether
\[
\forall\alpha<\omega_1\quad\exists\text{ a variety of groups }\mathcal V\quad h(\mathcal V)=\alpha,
\]
where \(\omega_1\) is the first uncountable ordinal, so the quantifier includes every finite ordinal and \(0\).

As contextual background, the supplied source also states that the cardinality of an ordinal height is finite, countably infinite, or equal to the cardinality of \(\omega_1\). This assertion is retained as supplied source context, without independent verification.
\subsection{Short English statement}
Does every countable ordinal occur as the supremum of the order types of all chains of proper subvarieties of some group variety that are well-ordered by increasing inclusion? The question is attributed to A. V. Kuznetsov (1980).
\subsection{Sources}
\begin{itemize}
\item Local UnsolvedMath registry, permanent ID \(700425\), exact source problem code \texttt{catalogue entry}, title \emph{Kourovka Notebook Problem 7.26}. The supplied record labels the problem \texttt{open}; this is a supplied status label, not an independently verified assessment.
\item \emph{The Kourovka Notebook (Unsolved Problems in Group Theory)}, original problem 7.26, 7th issue (1980), proposed by A. V. Kuznetsov. The supplied record identifies the referenced edition as the 21st issue edition (2026), \texttt{arXiv:1401.0300v48}.
\item The supplied collection description attributes editing of the Kourovka Notebook collection to E. I. Khukhro and V. D. Mazurov. No external source was consulted for this draft.
\end{itemize}
\end{document}
